\section{Cycle specialization and return multiplication}

Let $X=C_n$. Its fundamental group has one generator. The holonomy
representation is determined by one circuit receipt
\[
h=v_0v_1\cdots v_{n-1},
\]
and
\[
H_{\alpha}=\langle h\rangle.
\]
If $d=\ord(h)$, then $|H_{\alpha}|=d$. The component theorem gives
\[
\text{component count}=\frac{|R|}{d}.
\]
Every component contains $nd$ vertices. Because a regular lift of a cycle is
a disjoint union of cycles, every component is a cycle of length $nd$.

Thus Paper 6 is the rank-one case of the graph-wide theorem:
\[
\text{complete return length}
=
\text{visible circuit length}\times\text{receipt order}.
\]

The graph theorem adds something unavailable on one cycle. Several
independent circuit receipts can jointly generate all of $R$ even when no
single receipt has order $|R|$. For example, a noncyclic receipt group cannot
connect a lift over one cycle, but it can connect a lift over a graph whose
fundamental group has enough generators to surject onto that group.

\begin{corollary}[Generator-rank obstruction]
If $\widetilde X_{\alpha}$ is connected, then $R$ can be generated by at most
\[
\beta=|E|-|V|+1
\]
elements.
\end{corollary}

This is the first genuinely graph-wide constraint: the cycle rank of the
visible graph bounds the generator rank of any connected finite receipt
fiber.
